% GENERATED FILE. Edit canonical lessons, then run npm run generate:books.
% canonical-source-hash: fnv1a64:0e2cfc4f7756fa90
% canonical-lessons: MR-C76-navra, MR-C76-bayko, MR-C76-bal, MR-C76-manus, MR-C76-bai

\chapter{Husband, Wife, Baby, Person, Woman}
\label{ch:mr-husband-wife-baby-person-woman}
\hlchaptermodality{76}

\begin{chapteropening}
\textbf{By the end of this chapter:} \emph{I can say a husband, a wife, a baby, a person, a man and a woman.}

Complete the last lesson of chapter 76: I can say a husband, a wife, a baby, a person, a man and a woman.
\end{chapteropening}

\section[\texorpdfstring{\mr{नवरा} (navrā)}{नवरा (navrā)}]{\mr{नवरा} (navrā) — a husband}

\label{lesson:MR-C76-navra}



\begin{quote}
Before the new one: say the Marathi for an aunt, then the Marathi for an uncle.
\end{quote}

\subsection*{You'll want to know: \mr{नवरा}}
\textbf{\mr{नवरा}} — \emph{navrā} — \textquotedblleft{}a husband\textquotedblright{}.

\begin{cousinweb}
\textbf{\mr{नवरा}} (\emph{navrā}) is a tadbhava, a word inherited and reshaped by sound change. It comes from \emph{*navavaraka}, \textquotedblleft{}newly wed\textquotedblright{}, a form CDIAL reconstructs from its descendants (CDIAL 7013). It is \emph{nava}, \textquotedblleft{}new\textquotedblright{}, and \emph{vara}, \textquotedblleft{}bridegroom\textquotedblright{}.
\end{cousinweb}

\begin{grammarlens}[title={What you've built}]
One more word for this chapter.
\end{grammarlens}

\subsection*{Your turn}
\begin{itemize}
\raggedright
  \item \emph{Say it:} \emph{navrā}
  \item \emph{Say it:} \emph{navrā}, once more
  \item \emph{Say it:} say \emph{navrā}
  \item \emph{Recall:} say the Marathi for an aunt, then the Marathi for an uncle, then say \emph{navrā} again
\end{itemize}

\begin{tcolorbox}[breakable,colback=teal!4,colframe=teal!35!black,title={Before you move on}]
What does \mr{नवरा} mean? (\textquotedblleft{}A husband\textquotedblright{}.) Say it once more.
\end{tcolorbox}
\section[\texorpdfstring{\mr{बायको} (bāyko)}{बायको (bāyko)}]{\mr{बायको} (bāyko) — a wife}

\label{lesson:MR-C76-bayko}



\begin{quote}
Before the new one: say the Marathi for an uncle, then the Marathi for a husband.
\end{quote}

\subsection*{You'll want to know: \mr{बायको}}
\textbf{\mr{बायको}} — \emph{bāyko} — \textquotedblleft{}a wife\textquotedblright{}.

\begin{cousinweb}
The origin of \textbf{\mr{बायको}} (\emph{bāyko}) is not settled. Neither CDIAL nor DEDR, searched by form and by meaning, has an entry for it.
\end{cousinweb}

\begin{grammarlens}[title={What you've built}]
One more word for this chapter.
\end{grammarlens}

\subsection*{Your turn}
\begin{itemize}
\raggedright
  \item \emph{Say it:} \emph{bāyko}
  \item \emph{Say it:} \emph{bāyko}, once more
  \item \emph{Say it:} say \emph{bāyko}
  \item \emph{Recall:} say the Marathi for an uncle, then the Marathi for a husband, then say \emph{bāyko} again
\end{itemize}

\begin{tcolorbox}[breakable,colback=teal!4,colframe=teal!35!black,title={Before you move on}]
What does \mr{बायको} mean? (\textquotedblleft{}A wife\textquotedblright{}.) Say it once more.
\end{tcolorbox}
\section[\texorpdfstring{\mr{बाळ} (bāḷ)}{बाळ (bāḷ)}]{\mr{बाळ} (bāḷ) — a baby}

\label{lesson:MR-C76-bal}



\begin{quote}
Before the new one: say the Marathi for a husband, then the Marathi for a wife.
\end{quote}

\subsection*{You'll want to know: \mr{बाळ}}
\textbf{\mr{बाळ}} — \emph{bāḷ} — \textquotedblleft{}a baby\textquotedblright{}.

\begin{cousinweb}
\textbf{\mr{बाळ}} (\emph{bāḷ}) is a tadbhava, a word inherited from Sanskrit and reshaped by sound change. It comes from \emph{bāla}, \textquotedblleft{}young; boy\textquotedblright{} (CDIAL 9216). CDIAL lists the Marathi form \emph{bāḷ}, \textquotedblleft{}young\textquotedblright{}, and suggests the Sanskrit word perhaps belongs to a word group that gives several words for \textquotedblleft{}child\textquotedblright{}.
\end{cousinweb}

\begin{grammarlens}[title={What you've built}]
One more word for this chapter.
\end{grammarlens}

\subsection*{Your turn}
\begin{itemize}
\raggedright
  \item \emph{Say it:} \emph{bāḷ}
  \item \emph{Say it:} \emph{bāḷ}, once more
  \item \emph{Say it:} say \emph{bāḷ}
  \item \emph{Recall:} say the Marathi for a husband, then the Marathi for a wife, then say \emph{bāḷ} again
\end{itemize}

\begin{tcolorbox}[breakable,colback=teal!4,colframe=teal!35!black,title={Before you move on}]
What does \mr{बाळ} mean? (\textquotedblleft{}A baby\textquotedblright{}.) Say it once more.
\end{tcolorbox}
\section[\texorpdfstring{\mr{माणूस} (māṇūs)}{माणूस (māṇūs)}]{\mr{माणूस} (māṇūs) — a person, a man}

\label{lesson:MR-C76-manus}



\begin{quote}
Before the new one: say the Marathi for a wife, then the Marathi for a baby.
\end{quote}

\subsection*{You'll want to know: \mr{माणूस}}
\textbf{\mr{माणूस}} — \emph{māṇūs} — \textquotedblleft{}a person, a man\textquotedblright{}.

\begin{cousinweb}
\textbf{\mr{माणूस}} (\emph{māṇūs}) is a tadbhava, a word inherited from Sanskrit and reshaped by sound change. It comes from \emph{mānuṣa}, \textquotedblleft{}human; man\textquotedblright{} (CDIAL 10049). It is built on \emph{manu}, \textquotedblleft{}man\textquotedblright{}.
\end{cousinweb}

\begin{grammarlens}[title={What you've built}]
One more word for this chapter.
\end{grammarlens}

\subsection*{Your turn}
\begin{itemize}
\raggedright
  \item \emph{Say it:} \emph{māṇūs}
  \item \emph{Say it:} \emph{māṇūs}, once more
  \item \emph{Say it:} say \emph{māṇūs}
  \item \emph{Recall:} say the Marathi for a wife, then the Marathi for a baby, then say \emph{māṇūs} again
\end{itemize}

\begin{tcolorbox}[breakable,colback=teal!4,colframe=teal!35!black,title={Before you move on}]
What does \mr{माणूस} mean? (\textquotedblleft{}A person, a man\textquotedblright{}.) Say it once more.
\end{tcolorbox}
\section[\texorpdfstring{\mr{बाई} (bāī)}{बाई (bāī)}]{\mr{बाई} (bāī) — a woman}

\label{lesson:MR-C76-bai}



\begin{quote}
Before the new one: say the Marathi for a baby, then the Marathi for a person.
\end{quote}

\subsection*{You'll want to know: \mr{बाई}}
\textbf{\mr{बाई}} — \emph{bāī} — \textquotedblleft{}a woman\textquotedblright{}.

\begin{cousinweb}
The origin of \textbf{\mr{बाई}} (\emph{bāī}) is not settled. Neither CDIAL nor DEDR, searched by form and by meaning, has an entry for it.
\end{cousinweb}

\begin{grammarlens}[title={What you've built}]
One more word for this chapter.
\end{grammarlens}

\subsection*{Your turn}
\begin{itemize}
\raggedright
  \item \emph{Say it:} \emph{bāī}
  \item \emph{Say it:} \emph{bāī}, once more
  \item \emph{Say it:} say \emph{bāī}
  \item \emph{Recall:} say the Marathi for a baby, then the Marathi for a person, then say \emph{bāī} again
\end{itemize}

\begin{tcolorbox}[breakable,colback=teal!4,colframe=teal!35!black,title={Before you move on}]
What does \mr{बाई} mean? (\textquotedblleft{}A woman\textquotedblright{}.) Say it once more.
\end{tcolorbox}
